\documentclass[10pt,a4paper,sans]{moderncv}
\moderncvstyle{classic}
\moderncvcolor{blue}
\usepackage[utf8]{inputenc}
\usepackage[scale=0.82]{geometry}

% moderncv compose \today en anglais par défaut ; on redéfinit le format en français
% plutôt que de charger babel, qui modifierait la typographie de tout le document.
\renewcommand{\today}{\number\day\space\ifcase\month\or janvier\or février\or mars\or avril\or
mai\or juin\or juillet\or août\or septembre\or octobre\or novembre\or décembre\fi\space\number\year}

\name{Alexandro}{Disla}
\title{Lettre de Motivation}
\address{Rue Saint-Louis Jeanty \#14, Route de Frère}{Port-au-Prince, Haïti}{}
\phone[mobile]{(+509) 4148-3700}
\email{alexandrodisla@hotmail.com}
\extrainfo{GitHub: \url{https://github.com/AD0791}}

\begin{document}

\recipient{Ressources Humaines --- GOAL Haïti}{Pétion-Ville, Haïti\\Référence : Officier(ère) de redevabilité}
\date{\today}
\opening{Madame, Monsieur,}
\closing{Je vous prie d'agréer, Madame, Monsieur, l'expression de mes salutations distinguées,}
\enclosure[Ci-joint]{Curriculum Vitae, trois références professionnelles, attestation d'études}

\makelettertitle

Je vous adresse ma candidature au poste d'Officier de Redevabilité de GOAL Haïti. J'apporte cinq ans de suivi-évaluation en ONG en Haïti --- la santé et le VIH à la Caris Foundation International, l'eau et l'assainissement chez HANWASH, l'éducation chez Anseye Pou Ayiti --- précédés de huit ans de suivi d'indicateurs publics au Ministère de la Planification. Le créole haïtien et le français sont mes langues maternelles et je rédige aussi en anglais. Je réside à Route de Frère, à proximité de Pétion-Ville, et je suis disponible pour les visites de terrain comme pour le samedi.

Un mécanisme de retour d'information et de plaintes repose sur une promesse faite à chaque personne qui l'utilise : ce qu'elle a dit sera enregistré, transmis à qui peut agir, et lui reviendra sous forme de réponse. Tenir cette promesse exige une discipline que je pratique depuis cinq ans : une information juste dès la saisie, un suivi qui ne s'arrête qu'à la clôture, et des chiffres qui résistent à la vérification. À la Caris Foundation, j'ai établi et supervisé les protocoles d'assurance qualité des données d'un programme multi-sites, en remontant les chiffres jusqu'aux registres sources et en suivant chaque action corrective jusqu'à sa clôture plutôt que jusqu'à sa simple consignation. Les indicateurs que vous fixez au poste --- 90 \% des retours enregistrés sous 48 heures, 80 \% des dossiers clôturés dans les délais --- sont de ceux que je sais mesurer honnêtement et faire vivre dans un tableau de bord.

Le volet données et rapportage correspond à mon métier. J'ai administré des bases relationnelles, dédoublonné des registres de bénéficiaires, construit des tableaux de bord Power BI de suivi hebdomadaire et automatisé en Python et SQL la production de rapports d'activité, ce qui a réduit de 40 \% le temps de traitement manuel ; j'ai aussi produit et soumis des indicateurs destinés à l'USAID, avec les échéances et les contrôles de cohérence qu'impose un bailleur. Une analyse mensuelle des tendances n'a de valeur que si les catégories restent stables d'un mois à l'autre : c'est là que je serai le plus utile à l'équipe SERA. Quant à la confidentialité, j'ai travaillé trois ans sur les données d'un programme VIH, où une information mal gardée expose directement une personne. J'en retiens qu'une plainte sensible ne s'instruit pas au guichet : elle se réfère, avec discrétion, selon le protocole de sauvegarde.

Je préfère poser franchement l'écart principal. Je n'ai pas encore géré en titre un mécanisme de plaintes : mon expérience porte sur le suivi-évaluation et la gestion des données, pas sur la réception directe des plaintes. Ce qui s'y transfère est réel --- la conduite d'agents de terrain sur plusieurs sites, la formation d'équipes, le suivi de dossiers jusqu'à clôture et la maîtrise du système qui porte tout le reste. Ce qui s'apprend, les procédures CFM de GOAL et les normes AAP, je l'apprendrai comme j'ai appris tour à tour la santé, l'eau et l'éducation. De même, ma scolarité en économie appliquée au C.T.P.E.A est complétée et attestée par la pièce jointe, mais le diplôme reste en attente de la soutenance de mon mémoire de sortie.

L'approche de GOAL, fondée sur la participation communautaire, rejoint ce que je crois de la redevabilité : ce que disent les communautés doit modifier le programme, et non finir dans un classeur. Disponible dès octobre 2026, je reste à votre entière disposition pour un entretien.

\makeletterclosing

\end{document}
